\section{Appendix}

\subsection{Sample Output}
\label{appendix:sample_output}
\begin{minted}[fontsize=\small, frame=single, breaklines]{JSON}
{
  "claim_id": "C1",
  "claim_text": "X is a startup based in Germany.",
  "category": "other",
  "extraction_confidence": 0.95,
  "original_evidence": "'Startup Country': 'Germany'",
  "verification": {
    "verdict": "SUPPORTED",
    "certainty": 0.9,
    "reasoning": "Multiple independent third-party company/startup directories and profiles consistently list X's headquarters/location as Berlin, Germany. The company's own website also references Brandbrandnew GmbH and indicates operations ended, which aligns with the existence of a Germany-based startup entity.",
    "evidence_type": "consistent_sources",
    "source_count": 7,
    "sources": [
      "https://www.cbinsights.com/company/X",
      "https://startup-map.cologne/companies/X/t
      eam/alumni-founders",
      "https://www.startupnight.net/
      startups/2023/hX",
      "https://www.startbas
      e.com/organization/X/"
    ],
    "contradiction_details": null
  },
  "risk_assessment": {
    "risk_level": "low",
    "investor_alert": false,
    "requires_followup": false,
    "followup_questions": []
  }
}
\end{minted}


% =====================================


\subsection{Classification rules}
\label{appendix:classification_rules}

\begin{table*}[t]
\centering
\renewcommand{\arraystretch}{1.4}
\begin{tabular}{@{}p{0.2\linewidth} p{0.75\linewidth}@{}}
\toprule
\textbf{Annotation Field} & \textbf{Criteria and Rules} \\ 
\midrule

\multicolumn{2}{@{}l}{\textbf{1. Verdict Assignment ($L$)}} \\
\midrule
\texttt{SUPPORTED} & 
\textbf{Condition:} The retrieved evidence $E$ explicitly confirms or clearly entails the claim. \\
& \textbf{Rule:} Must be backed by at least one independent, high-reliability source (e.g., \texttt{Regulatory}, \texttt{News}, \texttt{Industry}). Statements found \textit{only} on the company's official channels (\texttt{Official}) cannot yield a \texttt{SUPPORTED} verdict on their own. \\

\texttt{CONTRADICTED} & 
\textbf{Condition:} The retrieved evidence $E$ explicitly refutes or clearly contradicts the claim. \\
& \textbf{Rule:} If evidence conflicts, the refuting evidence must be of equal or greater reliability than the supporting evidence to assign this label. \\

\texttt{INSUFFICIENT\_} & 
\textbf{Condition:} A definitive conclusion cannot be drawn from the available evidence. \\
\texttt{EVIDENCE} & \textbf{Rule:} Assign this label if: (i) no relevant evidence is found, (ii) evidence is entirely dependent or self-published (\texttt{Official}), (iii) evidence is vague or low-credibility (\texttt{Social}), or (iv) an unresolvable conflict exists among equally reliable sources. \\

\midrule
\multicolumn{2}{@{}l}{\textbf{2. Confidence Scoring ($s \in [0,1]$)}} \\
\midrule
\textbf{High Certainty} \newline ($0.8 \le s \le 1.0$) & 
\textbf{Condition:} The verdict is indisputable based on the available evidence. \\
& \textbf{Rule for \texttt{SUPPORTED}/\texttt{CONTRADICTED}:} Requires corroboration from multiple (at least two) independent, top-tier sources. \\

\textbf{Moderate Certainty} \newline ($0.5 < s < 0.8$) & 
\textbf{Condition:} The verdict is highly probable, but the evidence lacks perfect redundancy. \\
& \textbf{Rule:} Backed by exactly one independent, reliable source, or slightly ambiguous wording that leans strongly in one direction. \\

\textbf{Low Certainty} \newline ($0.0 \le s \le 0.5$) & 
\textbf{Condition:} The verdict is forced due to evidentiary weakness or severe ambiguity. \\
& \textbf{Rule:} This range is essentially reserved for \texttt{INSUFFICIENT\_EVIDENCE} verdicts, reflecting cautious abstention. \\
\end{tabular}
\caption{Annotation guidelines and rules for human claim evaluation. Evaluators assess each claim based on the retrieved evidence $E$ and assign a verdict $L$ along with a confidence score $s \in [0,1]$.}
\label{tab:annotation_rules}
\end{table*}

% ----------------------------------------------
\section{Primary System Prompts and Core Thresholds}
\label{app:prompts_thresholds}

VentureTruth's AI modules are driven by specialized prompts designed to enforce analytical rigor and conservative calibration. This appendix documents the core prompt intents together with the \emph{exact} thresholds and deterministic rules enforced throughout the pipeline.

\subsection{Claim Extraction}
\label{app:claim_extraction}

\textbf{Goal.} Extract verifiable marketing and factual claims from raw text.

\textbf{Filtering Rules.} The extraction prompt enforces strict exclusion of:
(i) funding goals and fundraising targets,
(ii) intentions and forward-looking statements (e.g., ``We plan to \dots''),
(iii) subjective opinions and value judgments, and
(iv) internal-only metadata (e.g., generic team sizes without external corroboration).
Compound factual statements (e.g., ``We partnered with X and grew Y\%'') are decomposed into distinct atomic claims.

\textbf{Limits and Decoding.} The extractor is configured to output between 3 and 20 claims per company (configurable), with a default maximum of 3 claims in the deployed configuration. Decoding is deterministic with temperature $T=0$.

\subsection{Evidence Retrieval (Search Manager)}
\label{app:evidence_retrieval}

\textbf{Goal.} Retrieve diverse, highly relevant public sources that may corroborate or refute extracted claims.

\textbf{Retrieval Settings.} Evidence retrieval uses Perplexity's \texttt{sonar-pro} API with a 60-second timeout and a maximum of 500 generated tokens per API call.

\textbf{Search Process.} The search prompt instructs the model to generate multi-angle queries (e.g., official, news, financial, and industry perspectives) to maximize coverage, targeting 5--10 distinct sources per claim.

\textbf{Quality Assessment and Retries.} Retrieval quality is evaluated on an ordinal scale from \textsc{FAILED} to \textsc{EXCELLENT}, where \textsc{EXCELLENT} corresponds to five or more highly relevant sources. If a retrieval attempt is rated \textsc{FAILED}, the system triggers exactly one automatic retry (maximum retries: 1) with forced query reformulation.

\subsection{Claim Verification}
\label{app:claim_verification}

\textbf{Goal.} Assess each claim against retrieved evidence and assign a verdict
$y \in \{\texttt{SUPPORTED}, \texttt{CONTRADICTED}, \texttt{INSUFFICIENT\_EVIDENCE}\}$.

\textbf{Certainty Output.} The verifier returns a calibrated certainty score $c \in [0,1]$ alongside the verdict.

\textbf{Strict Certainty Thresholds.} The system enforces the following hard constraints:
\begin{itemize}
    \item \textbf{No evidence found:} $c \le 0.3$.
    \item \textbf{Internal-only reference:} if evidence consists solely of internal material (e.g., an internal pitch deck) without public corroboration, the verdict is forced to \texttt{INSUFFICIENT\_EVIDENCE} and $c \le 0.4$.
    \item \textbf{Single source / interdependent sources:} $c \le 0.7$.
    \item \textbf{High confidence:} $c \ge 0.8$ requires at least two independent primary sources.
\end{itemize}

\textbf{Independence Definition.} Official website pages, press releases, and founder interviews are explicitly classified as \emph{dependent} sources and cannot be used alone to justify high-confidence support.

\subsection{Quality Assessment (Critic)}
\label{app:critic}

\textbf{Goal.} Critically evaluate verification iterations and produce feedback to guide iterative refinement.

\textbf{Settings.} The critic module is executed via \texttt{gpt-5.2-pro} with temperature $T=0.1$ for consistent analytical feedback.

\textbf{Structured Scoring.} The critic assigns continuous scores in $[0,1]$ along three dimensions:
\begin{enumerate}
    \item \emph{Reasoning Quality:} whether the reasoning is logically sound;
    \item \emph{Source Relevance:} whether sources are credible and independent;
    \item \emph{Certainty Calibration:} whether the reported certainty matches the enforced thresholds.
\end{enumerate}

\textbf{Output Categories and Actions.} Scores are mapped to five qualitative categories:
\textsc{EXCELLENT}, \textsc{GOOD}, \textsc{ACCEPTABLE}, \textsc{POOR}, and \textsc{CRITICAL}.
If the critic detects systematic failure modes, it outputs explicit action items (e.g., ``exclude press releases''), which are enforced in the subsequent iteration.

\subsection{Robustness Testing Variations}
\label{app:robustness}

\textbf{Goal.} Test verdict stability across different analytical personas and flag unstable verifications.

\textbf{Agreement Threshold.} A verification is flagged as unstable and may be forced to \texttt{INSUFFICIENT\_EVIDENCE} if the inter-variation agreement rate is below $0.67$ (e.g., fewer than two agreeing verdicts out of three prompt variations).

\textbf{Prompt Personas.} Robustness testing uses three prompt variants:
\begin{itemize}
    \item \textbf{Neutral (Default):} balanced verification under the standard rules above.
    \item \textbf{Evidence-Focused:} \texttt{SUPPORTED} requires $\ge 3$ independent high-quality sources; fewer than two sources default to \texttt{INSUFFICIENT\_EVIDENCE}.
    \item \textbf{Skeptical:} default stance is doubt; high confidence requires $\ge 4$ independent sources, and single-source verifications are penalized with $c \le 0.4$.
\end{itemize}

\subsection{Metadata Evaluation}
\label{app:metadata_eval}

Pipeline-extracted corporate metadata is validated against CRM ground truth (Salesforce) using token-level fuzzy string matching. A strict similarity threshold of $0.85$ is applied to tolerate benign lexical variations (e.g., ``Inc.'' vs.\ ``Incorporated''). Matching performance is reported using precision and recall.


% ================================

